\documentclass[10pt,a4paper]{amsart}
\usepackage{amsmath,amssymb,amsthm,graphicx}
\usepackage{verbatim}
\newenvironment{reference}{\comment}{\endcomment}
%\newenvironment{reference}{}{}
\newenvironment{slogan}{\comment}{\endcomment}
\newenvironment{history}{\comment}{\endcomment}
% For commutative diagrams we use Xy-pic
\usepackage[all]{xy}

% We use 2cell for 2-commutative diagrams.
\xyoption{2cell}
% Theorem environments.
%
\theoremstyle{plain}
\newtheorem{theorem}[subsection]{Theorem}
\newtheorem{proposition}[subsection]{Proposition}
\newtheorem{lemma}[subsection]{Lemma}

\theoremstyle{definition}
\newtheorem{definition}[subsection]{Definition}
\newtheorem{example}[subsection]{Example}
\newtheorem{exercise}[subsection]{Exercise}
\newtheorem{situation}[subsection]{Situation}

\theoremstyle{remark}
\newtheorem{remark}[subsection]{Remark}
\newtheorem{remarks}[subsection]{Remarks}

\numberwithin{equation}{subsection}

% Macros
%
\def\lim{\mathop{\mathrm{lim}}\nolimits}
\def\colim{\mathop{\mathrm{colim}}\nolimits}
\def\Spec{\mathop{\mathrm{Spec}}}
\def\Hom{\mathop{\mathrm{Hom}}\nolimits}
\def\Ext{\mathop{\mathrm{Ext}}\nolimits}
\def\SheafHom{\mathop{\mathcal{H}\!\mathit{om}}\nolimits}
\def\SheafExt{\mathop{\mathcal{E}\!\mathit{xt}}\nolimits}
\def\Sch{\mathit{Sch}}
\def\Mor{\mathop{\mathrm{Mor}}\nolimits}
\def\Ob{\mathop{\mathrm{Ob}}\nolimits}
\def\Sh{\mathop{\mathit{Sh}}\nolimits}
\def\NL{\mathop{N\!L}\nolimits}
\def\CH{\mathop{\mathrm{CH}}\nolimits}
\def\proetale{{pro\text{-}\acute{e}tale}}
\def\etale{{\acute{e}tale}}
\def\QCoh{\mathit{QCoh}}
\def\Ker{\mathop{\mathrm{Ker}}}
\def\Im{\mathop{\mathrm{Im}}}
\def\Coker{\mathop{\mathrm{Coker}}}
\def\Coim{\mathop{\mathrm{Coim}}}

% Boxtimes
%
\DeclareMathSymbol{\boxtimes}{\mathbin}{AMSa}{"02}

%
% Macros for moduli stacks/spaces
%
\def\QCohstack{\mathcal{QC}\!\mathit{oh}}
\def\Cohstack{\mathcal{C}\!\mathit{oh}}
\def\Spacesstack{\mathcal{S}\!\mathit{paces}}
\def\Quotfunctor{\mathrm{Quot}}
\def\Hilbfunctor{\mathrm{Hilb}}
\def\Curvesstack{\mathcal{C}\!\mathit{urves}}
\def\Polarizedstack{\mathcal{P}\!\mathit{olarized}}
\def\Complexesstack{\mathcal{C}\!\mathit{omplexes}}
% \Pic is the operator that assigns to X its picard group, usage \Pic(X)
% \Picardstack_{X/B} denotes the Picard stack of X over B
% \Picardfunctor_{X/B} denotes the Picard functor of X over B
\def\Pic{\mathop{\mathrm{Pic}}\nolimits}
\def\Picardstack{\mathcal{P}\!\mathit{ic}}
\def\Picardfunctor{\mathrm{Pic}}
\def\Deformationcategory{\mathcal{D}\!\mathit{ef}}
\usepackage{cleveref}
\crefname{originallemma}{original lemma}{original lemmas}
\Crefname{originallemma}{Original lemma}{Original lemmas}
\makeatletter
\newlabel{definition-gysin-homomorphism}{{02T8}{source}{}{sourcepointer-1}{}}
\newlabel{definition-gysin-homomorphism@cref}{{[originallemma][0][]02T8}{[1][0][]source}{}{}{}}
\makeatother
\crefname{originallemma}{original lemma}{original lemmas}
\Crefname{originallemma}{Original lemma}{Original lemmas}
\makeatletter
\newlabel{lemma-additivity-divisors-restricted}{{02QG}{source}{}{sourcepointer-2}{}}
\newlabel{lemma-additivity-divisors-restricted@cref}{{[originallemma][0][]02QG}{[1][0][]source}{}{}{}}
\makeatother
\crefname{originallemma}{original lemma}{original lemmas}
\Crefname{originallemma}{Original lemma}{Original lemmas}
\makeatletter
\newlabel{lemma-cycle-closed-coherent}{{02QY}{source}{}{sourcepointer-3}{}}
\newlabel{lemma-cycle-closed-coherent@cref}{{[originallemma][0][]02QY}{[1][0][]source}{}{}{}}
\makeatother
\crefname{originallemma}{original lemma}{original lemmas}
\Crefname{originallemma}{Original lemma}{Original lemmas}
\makeatletter
\newlabel{situation-setup}{{02QL}{source}{}{sourcepointer-4}{}}
\newlabel{situation-setup@cref}{{[originallemma][0][]02QL}{[1][0][]source}{}{}{}}
\makeatother
\title{The Gysin pullback of a cycle class and of a coherent sheaf class in the easy cases}
\author{The Stacks Project Authors}
\date{stacks-project:chow.tex:local-snapshot-20261012; GNU FDL 1.2 or later}
\begin{document}
\maketitle
\noindent\textit{Editorial note.} All delivered mathematical text is the author's own, copied from the pinned Stacks Project chapter file chow.tex; no statement or proof step was rewritten, reordered or removed.\par\medskip
\noindent\textit{Editorial note.} Preamble provenance: the pinned chapter file opens with \textbackslash{}input\{preamble\} and carries no preamble text of its own, so the project's single Stacks-wide preamble staged at sources/stacks-chow/preamble.tex is excerpted verbatim as three spans, lines 8-12 (the verbatim package together with the project's own reference, slogan and history display-environment definitions, including the author's commented-out older form of the reference environment), lines 14-18 (Xy-pic setup) and lines 154-217 (the native theorem declarations and mathematical macros). The lines 8-12 span is required because a counted statement or its surrounding text may use those display environments, which the project typesets as comments. Nothing was added, removed or redefined; the wrapper class and the amsmath/amssymb/amsthm/graphicx base packages are supplied by the assembler.\par\medskip
\noindent\textit{Editorial note.} Same-chapter pointers are located at their \textbackslash{}label lines inside the pinned chapter file; cross-chapter pointers are located in the project's own tag list staged in the snapshot (sources/stacks-chow/tags.txt), at the verbatim lines carrying their canonical tags.\par\medskip
\noindent\textit{Allowed standard prerequisites (editorial).}\begin{itemize}
\item The Gysin homomorphism attached to a triple of an invertible sheaf, a section and its zero scheme. Same chapter chow.tex, label definition-gysin-homomorphism (statement, proof line 5174). Fixes the pullback map i\textasciicircum{}* whose value on a cycle class and on a coherent sheaf class is being computed.
\item The cycle class of a closed subscheme and of its structure sheaf. Same chapter chow.tex, label lemma-cycle-closed-coherent, invoked at proof line 5198. Identifies the class of Z with that of its structure sheaf and the class of the intersection with the class of the structure sheaf of that intersection, which reduces part (1) to part (2).
\item Additivity of divisors for a restricted family. Same chapter chow.tex, label lemma-additivity-divisors-restricted, invoked at proof line 5233. Supplies the equality between length\_A(M/fM) and the sum over the minimal primes in the support of the multiplicity times the order of the local equation, which is the coefficient comparison that finishes the proof.
\end{itemize}\medskip
\ifdefined\hypertarget\hypertarget{sourcepointer-1}{}\fi
\noindent\textit{Source reference.} Original lemma 02T8: sources/stacks-chow/pinned-original-source.tex, source lines 5069--5069. Navigation pointer for \textbackslash{}ref\{definition-gysin-homomorphism\} used inside the counted span (line 5174); the argument that consumes it is retained in the span.\par\medskip
\ifdefined\hypertarget\hypertarget{sourcepointer-2}{}\fi
\noindent\textit{Source reference.} Original lemma 02QG: sources/stacks-chow/pinned-original-source.tex, source lines 373--373. Navigation pointer for \textbackslash{}ref\{lemma-additivity-divisors-restricted\} used inside the counted span (line 5233); the argument that consumes it is retained in the span.\par\medskip
\ifdefined\hypertarget\hypertarget{sourcepointer-3}{}\fi
\noindent\textit{Source reference.} Original lemma 02QY: sources/stacks-chow/pinned-original-source.tex, source lines 1776--1776. Navigation pointer for \textbackslash{}ref\{lemma-cycle-closed-coherent\} used inside the counted span (line 5198); the argument that consumes it is retained in the span.\par\medskip
\ifdefined\hypertarget\hypertarget{sourcepointer-4}{}\fi
\noindent\textit{Source reference.} Original lemma 02QL: sources/stacks-chow/pinned-original-source.tex, source lines 1372--1372. Navigation pointer for \textbackslash{}ref\{situation-setup\} used inside the counted span (line 5171); the argument that consumes it is retained in the span.\par\medskip
\begin{lemma}
\label{lemma-easy-gysin}
Let $(S, \delta)$ be as in Situation \ref{situation-setup}.
Let $X$ be locally of finite type over $S$.
Let $(\mathcal{L}, s, i : D \to X)$ be as in
Definition \ref{definition-gysin-homomorphism}.
\begin{enumerate}
\item Let $Z \subset X$ be a closed subscheme such
that $\dim_\delta(Z) \leq k + 1$ and such that
$D \cap Z$ is an effective Cartier divisor on $Z$. Then
$i^*[Z]_{k + 1} = [D \cap Z]_k$.
\item Let $\mathcal{F}$ be a coherent sheaf on $X$
such that $\dim_\delta(\text{Supp}(\mathcal{F})) \leq k + 1$ and
$s : \mathcal{F} \to \mathcal{F} \otimes_{\mathcal{O}_X} \mathcal{L}$
is injective. Then
$$
i^*[\mathcal{F}]_{k + 1} = [i^*\mathcal{F}]_k
$$
in $\CH_k(D)$.
\end{enumerate}
\end{lemma}

\begin{proof}
Assume $Z \subset X$ as in (1). Then set $\mathcal{F} = \mathcal{O}_Z$.
The assumption that $D \cap Z$ is an effective Cartier divisor is
equivalent to the assumption that
$s : \mathcal{F} \to \mathcal{F} \otimes_{\mathcal{O}_X} \mathcal{L}$
is injective. Moreover $[Z]_{k + 1} = [\mathcal{F}]_{k + 1}]$
and $[D \cap Z]_k = [\mathcal{O}_{D \cap Z}]_k = [i^*\mathcal{F}]_k$.
See Lemma \ref{lemma-cycle-closed-coherent}.
Hence part (1) follows from part (2).

\medskip\noindent
Write $[\mathcal{F}]_{k + 1} = \sum m_j[W_j]$ with $m_j > 0$
and pairwise distinct integral closed subschemes $W_j \subset X$
of $\delta$-dimension $k + 1$. The assumption that
$s : \mathcal{F} \to \mathcal{F} \otimes_{\mathcal{O}_X} \mathcal{L}$
is injective implies that $W_j \not \subset D$ for all $j$.
By definition we see that
$$
i^*[\mathcal{F}]_{k + 1} = \sum m_j [D \cap W_j]_k.
$$
We claim that
$$
\sum [D \cap W_j]_k = [i^*\mathcal{F}]_k
$$
as cycles.
Let $Z \subset D$ be an integral closed subscheme of $\delta$-dimension
$k$. Let $\xi \in Z$ be its generic point. Let $A = \mathcal{O}_{X, \xi}$.
Let $M = \mathcal{F}_\xi$. Let $f \in A$ be an element generating the
ideal of $D$, i.e., such that $\mathcal{O}_{D, \xi} = A/fA$.
By assumption $\dim(\text{Supp}(M)) = 1$,
the map $f : M \to M$ is injective, and
$\text{length}_A(M/fM) < \infty$. Moreover, $\text{length}_A(M/fM)$
is the coefficient of $[Z]$ in $[i^*\mathcal{F}]_k$. On the
other hand, let $\mathfrak q_1, \ldots, \mathfrak q_t$ be the minimal
primes in the support of $M$. Then
$$
\sum
\text{length}_{A_{\mathfrak q_i}}(M_{\mathfrak q_i})
\text{ord}_{A/\mathfrak q_i}(f)
$$
is the coefficient of $[Z]$ in $\sum [D \cap W_j]_k$.
Hence we see the equality by
Lemma \ref{lemma-additivity-divisors-restricted}.
\end{proof}

\end{document}
